\documentclass[11pt]{article}
\usepackage[margin=1in]{geometry}
\usepackage{amsmath,amssymb,amsthm,bm,xcolor,hyperref,booktabs,enumitem,graphicx}
\theoremstyle{plain}\newtheorem{theorem}{Theorem}\newtheorem{lemma}{Lemma}\newtheorem{proposition}{Proposition}
\theoremstyle{definition}\newtheorem{assumption}{Assumption}\newtheorem{remark}{Remark}
\newcommand{\E}{\mathbb{E}}\newcommand{\R}{\mathbb{R}}\newcommand{\N}{\mathcal{N}}\newcommand{\corr}{\mathrm{corr}}
\newcommand{\GP}{\mathcal{GP}}
\hypersetup{colorlinks=true,allcolors=blue!55!black}

\title{\textbf{Structure as Covariance: When a Graph Kernel Helps\\ Budget-Limited Active Retrieval}}
\author{Robert Richardson}
\date{\today}

\begin{document}\maketitle

\begin{abstract}
In multi-hop retrieval-augmented generation, relevance judgments are expensive LLM calls, so retrieval under a
judgment budget is a problem of Bayesian experimental design. We show that the corpus graph---which does
\emph{not} help as a ranking signal, because diffusion buries the answer---becomes load-bearing as the
\emph{covariance} of a latent relevance field with a semantic (retriever) mean: a judgment of one candidate
propagates belief to graph-connected candidates, surfacing a buried ``bridge'' passage that no single embedding
neighbour reveals. We make three contributions. \emph{(i)} A model---a Gaussian process with a semantic mean and a
graph covariance---together with the initially surprising finding that the graph kernel must be used in
\emph{correlation} (unit-diagonal) form, which roughly doubles its effect by equalizing hub variance so that
optimistic exploration reaches the propagating anchors. \emph{(ii)} An \emph{alignment law}: we prove, under a
stochastic block model, that the structural gain is zero when the graph is unaligned with the reasoning chain,
positive iff the graph is chain-\emph{assortative}, and monotone in a scalar alignment $p-q$ estimated in practice
by \emph{gold-connectivity}; the oracle-graph ceiling and the failure on adversarial corpora are the two ends of
this one axis. The same analysis explains why upper-confidence-bound acquisition beats value-of-information here:
the value of a judgment is the downstream surfacing it enables, an inherently non-myopic quantity. \emph{(iii)} An
empirical characterization under a real hop-aware judge across HotpotQA, 2WikiMultiHopQA, and MuSiQue: a
significant low-budget win on standard multi-hop (recall $+0.058$, nDCG@10 $+0.030$, chain-completion $+0.122$ at
one judgment), a partial rescue of adversarial MuSiQue by inferring the graph from the judge's own hop-reasoning
($+0.035$ at three hops, where every cheap surface graph is null), and an honest analysis of why the natural
adaptive gate does \emph{not} deploy---the per-query gain is a small distributional effect (signal-to-noise
$\approx 0.12$) rather than a per-query-predictable signal, so ``always use the graph'' is the correct, simple
policy. Effects are modest but principled, and the graph never significantly hurts.
\end{abstract}

\section{Introduction}
Retrieval-augmented generation increasingly relies on an LLM to \emph{judge} candidate passages, at a real
per-call cost. When the budget is a handful of judgments per query, retrieval becomes active learning: which
candidates should we spend judgments on, and how should a judgment reshape our beliefs about the rest? The hard
case is multi-hop: a question whose answer requires chaining facts across passages, where one gold passage---the
\emph{bridge}, supplying a linking entity rather than the answer itself---has low retriever similarity to the
query and sits buried among distractors. A bridge is exactly what an embedding retriever misses and what a
dependency structure could recover, if a judgment on a findable passage could \emph{propagate} to it.

\paragraph{Thesis: semantics sets the mean, structure sets the dependence.} We model latent relevance as a
Gaussian process whose \emph{mean} is a calibrated retriever score and whose \emph{covariance} is a corpus graph.
The graph is a poor ranking signal---using it to re-score \emph{buries} the bridge under diffusion---but an
excellent covariance: ``connected passages have correlated relevance,'' so conditioning on a judged anchor lifts
the posterior mean of a connected bridge. The same object that hurts as a mean helps as a dependence. This
resolves a tension in graph-based retrieval and, we argue, is the reason a Bayesian treatment earns its keep here:
the graph enters the \emph{decision/experimental-design} layer, not the ranking layer.

\paragraph{Contributions.}
\begin{enumerate}[leftmargin=*,itemsep=1pt,topsep=2pt]
\item \textbf{Model + the normalization finding} (\S\ref{sec:model}): a GP with a semantic mean and a graph
covariance, and the finding that the graph kernel must be in correlation (unit-diagonal) form---which doubles its
effect (chain-completion margin $+0.07\to+0.12$--$0.15$) by equalizing hub variance so optimistic exploration
reaches the propagating anchors.
\item \textbf{The alignment law} (\S\ref{sec:law}): a surfacing lemma and a stochastic-block-model theorem---the
structural gain is $0$ at zero alignment, positive iff the graph is chain-assortative ($p>q$), and monotone in
$p-q$; the oracle ceiling and the adversarial failure are the two ends of one axis, measured by gold-connectivity.
A corollary explains why UCB beats value-of-information (\S\ref{sec:acq}).
\item \textbf{An honest empirical characterization} (\S\ref{sec:exp}) under a real hop-aware judge: a significant
low-budget win on standard multi-hop; a partial rescue of adversarial MuSiQue via an LLM hop-assignment graph; and
an analysis of why per-query adaptive routing does not deploy, yielding the simple ``always use the graph'' policy.
\end{enumerate}
We do not claim large gains: effects are modest. We claim a principled account---\emph{when} structure helps and
\emph{why}---and a method that is safe (never significantly hurts) and simple.

\section{Related work}\label{sec:rel}
\emph{Active / Bayesian retrieval.} Our nearest neighbour is BAGEL, a query-specific GP-UCB over passage
embeddings; our data refute the embedding kernel here and replace it with a graph covariance, and we add a soft
(noisy-likelihood) update that prevents a GP from confidently propagating a wrong LLM label---BAGEL's two named
limitations. \emph{Graph retrieval.} GraphRAG and HopRAG use graph structure for retrieval; HopRAG builds logical
pseudo-query edges precisely because surface similarity is insufficient, consonant with our alignment law. We use
the graph as a covariance, not a ranking signal, and characterize when it helps. \emph{GMRF/CAR priors} (Rue \&
Held) supply the kernel; \emph{Bayesian optimization / value of information} the acquisition; \emph{set-wise
retrieval} (SETR) the completion metric. Our novelty is the mean/dependence decomposition, the exploration
mechanism, and the alignment regime---not any single primitive.

\section{Model: structure as covariance}\label{sec:model}
Per query, a first-stage retriever returns $N$ candidates with normalized embeddings $v_i$ and retriever scores;
a calibrated logistic map gives a prior mean $m_i\in[0,1]$ (cross-fit cosine$\to$gold). We place a GP on latent
relevance $f\sim\GP(m,K)$ and, when candidate $j$ is judged, observe $y_j\in\{0,\tfrac12,1\}$ (a graded hop-aware
label) under Gaussian noise $\sigma^2$. Ranking is by posterior mean; recall@$k$, chain-completion, and nDCG@10
are measured against gold.

\paragraph{Two kernels.} The embedding kernel $K^E_{ij}=\exp(-(1-v_i^\top v_j)/\ell)$ encodes ``similar
embeddings, similar relevance''---but the bridge is dissimilar to everything, so $K^E$ cannot reach it. The graph
kernel is a GMRF/CAR covariance $K^G=(I+\lambda \mathcal L)^{-1}$ on a corpus graph $A$ (Laplacian
$\mathcal L=D-A$), encoding ``connected passages, correlated relevance.''

\paragraph{The load-bearing detail: correlation form.} We find the graph kernel must be normalized to unit
diagonal, $\widetilde K^G=\mathrm{diag}(K^G)^{-1/2}K^G\,\mathrm{diag}(K^G)^{-1/2}$, before use. Unnormalized, high-degree
hubs have \emph{lower} prior variance, so an upper-confidence-bound rule under-explores them---but the hubs are
exactly the anchors whose judgment propagates to bridges. Equalizing variance lets exploration reach the hubs; the
propagation then fires. Empirically this $1.8\times$ the graph$-$cosine chain-completion margin ($+0.068\to+0.122$
at one judgment) and, unlike the raw kernel (which converges to null by three judgments), keeps the win
significant at every budget. The finding is a small but real methodological point and is, to our knowledge, not
standard in GP-active-retrieval practice.

\paragraph{The soft update.} Conditioning on judged set $J$, $\mu = m + K_{:J}(K_{JJ}+\sigma^2 I)^{-1}(y_J-m_J)$.
We rank by $\mu$ (soft) rather than hard-pinning judged labels; with an imperfect judge, hard pinning of a wrong
label is catastrophic (a mislabeled gold is sunk), while the soft update degrades gracefully.

\section{Acquisition: exploration, not information}\label{sec:acq}
We select judgments by UCB, $\arg\max_i \mu_i+\beta\,\mathrm{sd}_i$. We tested principled alternatives---one-step
value of information for chain completion ($1-\prod_j(1-p_j)$), two-step lookahead, pure information gain (maximum
variance reduction), and pure exploration (maximum variance)---and \emph{all lose} to UCB on the graph kernel
(chain-completion $-0.08$ to $-0.11$ at one judgment; two-step lookahead does not help). The reason, made precise
by the alignment analysis below, is that the value of judging a high-prior anchor is not its own label but the
\emph{downstream} surfacing of a connected bridge---an inherently multi-stage, non-myopic quantity that
single-step VOI scores at nearly zero. UCB captures it because the mean term focuses exploration on the high-prior
anchors along which belief propagates; kernel and acquisition are entangled, and the graph \emph{needs} UCB.

\section{The alignment law}\label{sec:law}
When does the graph help? We give an exact surfacing condition and lift it to a statistical law.
\begin{assumption}[Buried bridge]\label{ass:bridge}
Gold $=\{a,b\}$ with $m_a=\max_i m_i$ (findable anchor) and $m_b<\min_{d}m_d$ (bridge buried below distractors
$D$); judging is informative, $y_a>m_a$.
\end{assumption}
\begin{lemma}[Surfacing]\label{lem:surface}
With unit prior variances the first UCB pick is $a$; conditioning gives $\mu_i=m_i+\frac{K_{ia}}{1+\sigma^2}
(y_a-m_a)$, and the bridge enters the top-$k$ iff the \emph{kernel differential} exceeds the burial threshold,
\[
   K_{ba}-\max_{d\in D}K_{da}\ >\ \tau:=\frac{(1+\sigma^2)(\max_d m_d-m_b)}{y_a-m_a}.
\]
\end{lemma}
The mechanism is isolated: a judgment surfaces the bridge only through a kernel that couples it to the findable
anchor \emph{more} than it couples distractors. The embedding kernel has this differential $\le 0$ for a buried
bridge, so it never surfaces it; the graph kernel can.

\begin{theorem}[Structure helps iff chain-assortative]\label{thm:align}
Let the graph be drawn from a stochastic block model with within-chain edge probability $p$ and off-chain $q$
(alignment $p-q$). Then \emph{(i)} for the one-hop kernel $\corr(I+\beta A)$, $\E[K_{ba}-\max_d K_{da}]$ is
strictly increasing in $p-q$ with leading term $\beta(p-q)$ and $0$ at $p=q$ (monotone for the diffusion kernel);
\emph{(ii)} since the embedding kernel never surfaces the bridge, the expected structural gain
$\Delta(p,q)=\Pr[\text{surface}\mid K^G]-\Pr[\text{surface}\mid K^E]$ is $0$ at $p=q$, positive iff $p>q$, and
monotone increasing in $p-q$; \emph{(iii)} $q\to0,p\to1$ (gold clique) is the oracle-graph ceiling, and $p\approx
q$ (edges independent of the chain---surface-similar distractors) gives $\Delta\approx0$.
\end{theorem}
Proof in Appendix~\ref{app:proof}. We verify it on a synthetic SBM instance (Figure~\ref{fig:align}A): the graph
advantage rises monotonically from $+0.007$ at $p=q$ to $+0.115$ at $p-q=0.95$, essentially $0$ at $p=q$, while the
embedding kernel stays at chance. The real data (Figure~\ref{fig:align}B) trace the same law, from aligned
(chained) through neutral (comparison) to the adversarial MuSiQue boundary.
\begin{figure}[t]\centering
\includegraphics[width=\linewidth]{fig_alignment.pdf}
\caption{The alignment law. \textbf{(A)} Synthetic SBM: the graph kernel's gain rises with graph--chain alignment
$p-q$ and is $\approx 0$ at $p=q$; the embedding kernel is at chance throughout. \textbf{(B)} Real graph$-$cosine
margins under a real hop-aware judge: positive on chained (aligned), neutral on comparison, null for cheap graphs
on adversarial MuSiQue, partially recovered by an LLM hop-assignment graph, up to the oracle ceiling. Bars are
$95\%$ CIs.}\label{fig:align}
\end{figure}
\begin{remark}[Measurable alignment]
In practice $p-q$ is estimated by \emph{gold-connectivity} (are gold passages connected to each other more than to
distractors). Theorem~\ref{thm:align} explains why this single scalar predicts the gain, why the oracle ceiling
and the adversarial failure are one axis, and why the effect saturates.
\end{remark}

\section{Experiments}\label{sec:exp}
\paragraph{Setup.} Pools of $N=100$ retrieved from a full-dev corpus (so the bridge is genuinely buried and the
prior is weak); a graded $0$--$2$ hop-aware LLM judge (gpt-4o-mini) that explicitly credits linking passages; the
soft GP with $\sigma^2$ tuned to judge reliability. Baselines: no-judge prior, passive verify-top-$B$,
embedding-kernel GP (BAGEL-lite), and ceilings (oracle judge, oracle graph). Metrics: recall@$k$, chain-completion
($k=|\text{gold}|$), nDCG@10; paired bootstrap 95\% CIs; query-grouped splits. Datasets: HotpotQA,
2WikiMultiHopQA, MuSiQue. A naive judge (``does this answer the question?'') is bridge-blind (gold recall $0.35$);
the hop-aware prompt raises gold recall to $0.84$, isolating the prompt fix from model strength.

\paragraph{Main result: a significant low-budget win on standard multi-hop.} On chained HotpotQA/2Wiki
($n{=}600$), with the correlation-form kernel the graph beats the embedding kernel across all three metrics and all
budgets (Table~\ref{tab:main}): at one judgment, recall $+0.058\,[.040,.076]$, nDCG@10 $+0.030\,[.020,.041]$,
chain-completion $+0.122\,[.088,.155]$; significant through $B{=}3$. The graph kernel---not ``any GP''---drives it:
graph-GP $>$ cosine-GP $>$ passive.
\begin{table}[t]\centering\small
\begin{tabular}{lccc}
\toprule
graph$-$cosine (chained, real judge) & recall@$k$ & nDCG@10 & completion \\ \midrule
$B=1$ & $+0.058\,[.040,.076]$ & $+0.030\,[.020,.041]$ & $+0.122\,[.088,.155]$ \\
$B=2$ & $+0.071\,[.052,.090]$ & $+0.032\,[.021,.045]$ & $+0.150\,[.115,.185]$ \\
$B=3$ & $+0.063\,[.044,.082]$ & $+0.033\,[.022,.045]$ & $+0.127\,[.093,.162]$ \\
\bottomrule
\end{tabular}
\caption{Low-budget graph-kernel win (correlation-form kernel), significant at all budgets. The raw kernel is
weaker and converges to null by $B=3$ (\S\ref{sec:model}, normalization).}\label{tab:main}
\end{table}

\paragraph{The alignment law, from both sides.} We judged the \emph{comparison} (independent-evidence) questions
as a control, giving a mixed distribution ($300$ chained $+$ $300$ comparison). The graph helps chained
($+0.058\,[.040,.076]$ at $B{=}1$) and does \emph{not} help comparison (neutral to slightly negative,
$\le 0$)---the boundary the theorem predicts, with the mechanism visible: comparison questions have a strong prior
(passive recall $\sim0.8$; both entities directly findable, no bridge to propagate along). The oracle-graph
ceiling is $+0.07$--$0.09$, and gold-connectivity tracks the per-regime gain.

\paragraph{Adversarial corpora: MuSiQue, and a partial rescue.} MuSiQue's distractors are made surface-similar to
golds, so cheap graphs (title, entity-overlap, cosine-decomposition) have $p\approx q$ and are null
(graph$-$cosine $+0.008\,[\text{-}.02,.03]$ at three hops), while an oracle gold-clique graph recovers
$+0.087\,[.05,.13]$---exactly the theorem's two ends. Following its prescription (raise $p-q$), we build the graph
from the judge's \emph{own} reasoning: decompose the question into sub-questions and ask the judge which
sub-question each pool passage answers, then connect passages on different hops. This recovers a significant
$+0.035\,[.01,.06]$ at three hops---$\sim4\times$ the cosine graph and $\sim40\%$ of the oracle headroom, robust
across graph constructions---where every surface graph is null. It is a partial rescue, concentrated in the deep
multi-hop regime; a proof-of-concept that inferring structure from the judge's reasoning is the route to
adversarial corpora.

\paragraph{Why not adapt per query? An honest negative.} A natural extension is a gold-free gate that decides
per query whether to use the graph. It does not deploy, and the reason is instructive and robust to the kernel
choice. The per-query graph advantage is a \emph{small mean effect swamped by variance} (signal-to-noise
$\approx0.1$--$0.2$) and is nearly unpredictable from gold-free features \emph{even under the oracle}: ridge
$R^2\approx0.03$--$0.05$, and only marginally better with theory-motivated alignment/reachability features. Two
mechanistic fixes we tried failed: the graph does \emph{not} amplify judge errors (advantage correlates $-0.17$
with anchor reliability), and confidence-gating the propagation to high-confidence labels \emph{hurts} ($-0.030$;
the ``related'' grade-$1$ labels carry useful signal). On the mixed distribution the learned gate only \emph{ties}
always-graph ($+0.003$, n.s.) while the per-query oracle has a real $+0.027$ headroom that is simply not
gold-free-reachable---and this holds under the correlation-form kernel, even though there the graph slightly
\emph{penalizes} comparison queries. The structural gain is thus a \emph{distributional} property---a robust
average in the deep multi-hop regime (the alignment law)---not a per-query-predictable signal. The correct policy
is to use the graph in the multi-hop regime it targets; per-query adaptivity is neither feasible nor needed.

\section{Discussion and limitations}\label{sec:disc}
Effects are modest and low-budget; the win is over a re-implemented BAGEL; two of three benchmarks are standard
multi-hop. The contribution is the account: a mean/dependence decomposition, a proved alignment law that unifies
the oracle ceiling and the adversarial boundary, an exploration mechanism, and an honest map of where the method
helps (deep multi-hop with a chain-assortative graph), is neutral (independent evidence), and is only partially
recoverable (adversarial corpora, via inferred structure). The deepest open direction is \emph{adaptive structure
learning}: inferring the reasoning-chain graph from the first judgments---of which the MuSiQue hop-assignment
result is a first instance---which would raise $p-q$ within the budget and is the most plausible route to helping
on adversarial corpora without a gold-derived graph.

\appendix
\section{Proof of Theorem~\ref{thm:align}}\label{app:proof}
\emph{(i)} After unit-diagonal normalization the one-hop kernel has $K_{ij}=\beta A_{ij}$ for $i\neq j$, so
$\E[K_{ba}]=\beta p$, and the anchor--distractor edges being i.i.d.\ $\mathrm{Bernoulli}(q)$ gives
$\E[\max_d K_{da}]=\beta(1-(1-q)^{|D|})$, whence $\E[K_{ba}-\max_d K_{da}]=\beta(p-1+(1-q)^{|D|})$: strictly
increasing in $p$, decreasing in $q$, zero at $p=q$ to first order, $\le0$ for $p\le q$. The diffusion kernel adds
$\rho^\ell$-weighted path counts, each increasing in $p$ relative to $q$, preserving sign and monotonicity.
\emph{(ii)} By Lemma~\ref{lem:surface} the bridge surfaces under $K$ iff its differential exceeds $\tau>0$; the
embedding differential is $\le0<\tau$, so the second term of $\Delta$ vanishes, and the first,
$\Pr[K^G_{ba}-\max_d K^G_{da}>\tau]$, increases in $p-q$ by (i), is $0$ when the differential concentrates at $0$
($p=q$), and positive once $p>q$. \emph{(iii)} At $q{=}0,p{=}1$ the differential is maximal; at $p=q$ it is $0$.
\hfill$\square$

\smallskip\noindent\footnotesize\emph{Reproducibility.} All experiments use cached judgments; scripts and the full
experimental log accompany the submission. Total LLM spend for all runs reported here is under \$10.
\end{document}
